\BeleskeID{08-s053}
% element: 08-s053:e03
% element: 08-s053:notes
Polazi se od gotovog rešenja umesto ponovnog pisanja i rešavanja diferencijalne jednačine. Najpre se određuje $\tau$ kao proizvod $C$ i otpornosti koju kondenzator vidi uz kratko spojene nezavisne naponske izvore i uklonjene nezavisne strujne izvore, odnosno otvorene veze. Zavisni izvori ostaju.

Početna vrednost tražene veličine $p(t_0^+)$ određuje se iz realnog kola uz saznanje da se napon kondenzatora ne može promeniti trenutno. To ne znači automatski da svaka druga veličina ima kontinuitet.

Krajnja vrednost $p(\infty)$ dobija se po završetku prelaznih pojava. Tada nema promene napona kondenzatora, pa je njegova struja nula. Sve tri veličine zatim se unose u istu formulu odziva. Konačno stanje bira se za pobudu koja važi u razmatranom intervalu.
